\documentclass{article}
\usepackage{amsmath}
\usepackage{amsfonts}
\usepackage{amssymb}
\usepackage{graphicx}
\usepackage{matlab-prettifier}
\usepackage{IEEEtrantools}

\setlength\parindent{0pt}

\begin{document}


\title{Systems of ODEs}
\author{AE 2440}
\date{}
\maketitle


Numerical solvers like \texttt{ode45} in MATLAB expect your system (or model) to be expressed as a first-order system:

\begin{lstlisting}[style=Matlab-editor]
    dxdt = f(t, x)
\end{lstlisting}
where \( x \) is a vector of state variables, and \( f(t, x) \) is a function that describes the dynamics of the system.
This document provides a brief overview of how to convert higher-order ordinary differential equations (ODEs) into a system of first-order ODEs, which can then be solved using numerical methods like \texttt{ode45}.


\section*{General Form of a System of ODEs}

A system of ordinary differential equations can be written in vector form as:

\[
\frac{d\mathbf{x}}{dt} = \mathbf{f}(t, \mathbf{x}(t))
\]

where:
\begin{itemize}
  \item $\mathbf{x}(t) \in \mathbb{R}^n$ is the vector of unknown functions (state variables),
  \item $t$ is the independent variable (often time),
  \item $\mathbf{f}: \mathbb{R} \times \mathbb{R}^n \to \mathbb{R}^n$ is a vector-valued function defining the dynamics.
\end{itemize}

\subsection*{Example}

For a 2-dimensional system:

\[
\begin{cases}
\frac{dx}{dt} = f_1(t, x, y) \\
\frac{dy}{dt} = f_2(t, x, y)
\end{cases}
\]

Note that \textbf{the first derivative of each state variable is expressed in terms of only the state variables themselves and possibly time}. This is a key feature of systems of ODEs, allowing for the representation of complex dynamics in a structured manner.  

\section*{N-th order ODE as N first-order ODEs}

Any \( n \)-th order ODE can always be reformulated as a system of \( n \) first-order ODEs. This approach is fundamental in mathematical modeling and numerical methods, as it allows higher-order equations to be solved using methods that are designed for systems of first-order equations, like MATLAB’s \texttt{ode45}.

Here’s a brief outline of the process:

\begin{enumerate}
    \item \textbf{Define State Variables}: For an \( n \)-th order ODE, introduce a set of new variables to represent the derivatives up to the \( (n-1) \)-th order. If we start with an \( n \)-th order equation:
    \[
    y^{(n)} = f(t, y, y', \dots, y^{(n-1)})
    \]
    we define \( y_1 = y \), \( y_2 = y' \), up to \( y_n = y^{(n-1)} \).
    
    \item \textbf{Rewrite the System}: Now express the original equation in terms of these state variables, yielding a system of first-order ODEs:
    \[
    \begin{cases}
        \dot{y}_1 = y_2, \\
        \dot{y}_2 = y_3, \\
        \vdots \\
        \dot{y}_{n-1} = y_n, \\
        \dot{y}_n = f(t, y_1, y_2, \dots, y_n).
    \end{cases}
    \]
    This converts the single \( n \)-th order ODE into a system of \( n \) first-order ODEs, which can then be solved with methods such as \texttt{ode45}.
\end{enumerate}

This transformation is universally applicable because it leverages the fundamental property that any \( n \)-th order derivative can be expressed as a sequence of first-order derivatives, allowing the problem to be represented in state-space form.

\section*{Example: Mass-Spring-Damper System, AKA Damped Harmonic Oscillator}

Consider the second-order ODE of a damped harmonic oscillator:
\[
m \ddot{y} + c \dot{y} + k y = 0
\]
with initial conditions:
\begin{IEEEeqnarray*}{rCl}
y(t=0) &= & 1.0 \: \text{m}\\
\dot{y}(t=0) &=& 0 \: \text{m/s}
\end{IEEEeqnarray*}
where:
\begin{itemize}
    \item \( y \) is the displacement,
    \item \( m \) is the mass,
    \item \( c \) is the damping coefficient, and
    \item \( k \) is the spring constant.
\end{itemize}

We can rewrite this in the standard form:
\[
\ddot{y} = -\frac{c}{m} \dot{y} - \frac{k}{m} y
\]

\subsection*{Define the State Variables}

To convert this into a system of first-order ODEs, we define two new state variables:
\[
y_1 = y \quad \text{(position) and} \quad y_2 = v_y =  \dot{y} \quad \text(velocity).
\]

\subsection*{Define the State Vector}

It is important to explicity define the convention you will use for the order in the state vector, $\mathbf{x}(t)$.  There is not right or wrong order to use, but you will need to be consistent when implementing the problem in software, so it is important to write down your intentions.

We could choose this order:
\[
\mathbf{x}(t) =
\begin{Bmatrix} y \\ v_y \end{Bmatrix} =
\begin{Bmatrix} y_1 \\ y_2 \end{Bmatrix}.
\]

(We could also, equivalently choose the opposite order:
\[
\mathbf{y}(t) =
\begin{Bmatrix}  v_y \\ y \end{Bmatrix} =
\begin{Bmatrix} y_2 \\ y_1 \end{Bmatrix}.
\]
We just need to be consistent with either convention.)


\subsection*{Rewrite as a System of First-Order ODEs}

This allows us to write the system as:
\[
\begin{cases}
\dot{y} = v_y \\
\dot{v}_y = -\frac{c}{m} v_y - \frac{k}{m} y
\end{cases}
\]

This system is now ready for solving with MATLAB’s \texttt{ode45}.

\subsection*{Step 3: Implement in MATLAB}

To solve this system, define parameters \( m \), \( c \), and \( k \), as well as initial conditions. Use the following MATLAB code:

\begin{lstlisting}[style=Matlab-editor]
% Define parameters
m = 1;       % mass
c = 0.2;     % damping coefficient
k = 5;       % spring constant

% Define initial conditions [y(0); y'(0)]
y0 = [1; 0];  % initial displacement = 1, initial velocity = 0

% Time span
tspan = [0 10];

% Define the system of equations as anonymous function
dydt = @(t, y) [y(2);
                -c/m * y(2) - k/m * y(1)];

% Solve the system with ode45
[t, y] = ode45(dydt, tspan, y0);

% Plot the results
plot(t, y(:,1))
xlabel('Time (s)')
ylabel('Displacement y(t)')
title('Damped Harmonic Oscillator')
grid on
\end{lstlisting}

\subsection*{Explanation of the Code}

\begin{enumerate}
    \item \textbf{Define Parameters}: Set values for \( m \), \( c \), and \( k \).
    \item \textbf{Define Initial Conditions}: Start with initial displacement \( y(0) = 1 \) and initial velocity \( y'(0) = 0 \).
    \item \textbf{Set Up System}: Use an anonymous function to represent the system of first-order ODEs.
    \item \textbf{Call \texttt{ode45}}: Solve the system over the time interval \([0, 10]\).
    \item \textbf{Plot Results}: Plot \( y(t) \), the displacement, to observe the damped oscillations.
\end{enumerate}

By following these steps, you can solve and visualize the behavior of the damped harmonic oscillator using MATLAB’s \texttt{ode45}.

\end{document}
